% ch5_answers.tex -- answers to the Quick check questions of Module 8, Chapter 5.
\begin{chapteranswers}{5}{Deep Inside the Atom}
  % ---- 5.1 ----
  \item A fundamental particle has no smaller parts (no inner structure), as far as
        experiments can tell. Yes: the electron is fundamental. Experiments find no size
        and no parts inside it.
  \item $\lambda = \dfrac{h}{mv} = \dfrac{\qty{6.626e-34}{J.s}}{\qty{8.0e-18}{kg.m.s^{-1}}}
        = \qty{8.28e-17}{m}$ (the momentum $mv$ is given). Compare:
        $\dfrac{\qty{1.7e-15}{m}}{\qty{8.28e-17}{m}} = 21$. Answer:
        \final{$\lambda = \qty{8.3e-17}{m}$}. The wavelength is about 21 times
        \emph{smaller} than a proton, so yes, these electrons can show structure inside it.
  \item In beta-minus decay a neutron changes into a proton:
        $\mathrm{n} \to \mathrm{p} + \mathrm{e}^- + \overline{\nu}_\mathrm{e}$. In the
        quark model this is udd $\to$ uud: only one down quark changes into an up quark,
        through a $\mathrm{W}^-$ boson. The other two quarks stay the same. The neutron
        turns into a proton when just one of its parts changes. This is easy to explain if
        the neutron has parts. So beta decay supports the idea that the neutron is not
        fundamental. (It is not proof on its own, because the muon is fundamental and also
        decays.)
  \item Any two, each with its evidence. Positron: a cloud-chamber track with an
        electron's mass that bent the ``positive'' way; a lead plate showed which way it
        moved (Anderson, 1932). Neutrino: beta electrons have a spread of energies, so
        some energy seemed to go missing (Pauli, 1930); it was then detected near a
        nuclear reactor (1956). Muon: cosmic-ray tracks of a particle with charge $-e$ but
        207 times the electron's mass (1936). Higgs boson: a bump at
        \qty{125}{GeV}$/c^2$ in LHC data (2012).
  % ---- 5.2 ----
  \item Each antiquark has the opposite sign of its quark:
        $Q = -\tfrac23 - \tfrac23 + \tfrac13 = \dfrac{-2 - 2 + 1}{3} = -1$.
        Answer: \final{$Q = -1e = \qty{-1.602e-19}{C}$}.
  \item Muon: lepton. $\pi^-$: meson. Proton: baryon. Photon: force carrier (exchange
        particle of the EM force). Electron neutrino: lepton. Fundamental: the muon, the
        photon and the electron neutrino. The $\pi^-$ and the proton are made of quarks.
  \item First-generation quarks are u ($+\tfrac23$) and d ($-\tfrac13$), and a baryon has
        three quarks. Try each mix: ddd $= -1$, udd $= 0$, uud $= +1$,
        uuu $= \tfrac23 + \tfrac23 + \tfrac23 = \tfrac63 = +2$. Only uuu works.
        Answer: \final{uuu} (this baryon is the $\Delta^{++}$).
  \item Gravity (graviton, never found) $<$ weak force ($\mathrm{W}^+$, $\mathrm{W}^-$,
        $\mathrm{Z}^0$) $<$ electromagnetic force (photon) $<$ strong force (gluons).
  \item An up quark changes into a down quark: $\mathrm{u} \to \mathrm{d} + \mathrm{W}^+$,
        then $\mathrm{W}^+ \to \mathrm{e}^+ + \nu_\mathrm{e}$. Overall, inside a nucleus,
        $\mathrm{p}\ (\mathrm{uud}) \to \mathrm{n}\ (\mathrm{udd}) + \mathrm{e}^+ +
        \nu_\mathrm{e}$. (A free proton cannot do this: it is lighter than a neutron. The
        nucleus supplies the energy.) The boson is the \final{$\mathrm{W}^+$}, a carrier
        of the weak force.
  % ---- 5.3 ----
  \item $W = qV = 2 \times \qty{1.602e-19}{C} \times \qty{5.0e4}{V} = \qty{1.602e-14}{J}$.
        In electronvolts: a charge of $2e$ through \qty{50}{kV} gains
        $2 \times \qty{50}{keV} = \qty{100}{keV}$. Answer:
        \final{$\qty{100}{keV} = \qty{1.6e-14}{J}$}.
  \item $f = \dfrac{1}{T} = \dfrac{qB}{2\pi m} = \dfrac{\qty{1.602e-19}{C} \times
        \qty{0.80}{T}}{2\pi \times \qty{1.673e-27}{kg}} = \qty{1.219e7}{Hz}$. Answer:
        \final{$f = \qty{1.2e7}{Hz}$}. $f$ is proportional to $B$, so doubling $B$
        doubles the frequency, to \qty{2.4e7}{Hz}.
  \item The magnetic force is always at \ang{90} to the velocity. So it changes only the
        direction, not the speed. But the cyclotron still needs it, for three reasons.
        First, it bends the protons into half circles, so they return to the gap again and
        again and get many pushes. Second, it sets the period $T = \dfrac{2\pi m}{qB}$, so
        a fixed-frequency AC stays in step. Third, with the dee radius $R$, it sets the
        top speed, $v = \dfrac{qBR}{m}$.
  \item The new mass is a tau plus an antitau (same mass):
        $E = 2m_\tau c^2 = 2 \times \qty{3.167e-27}{kg} \times (\qty{3.00e8}{m.s^{-1}})^2
        = \qty{5.701e-10}{J}$. In eV:
        $\dfrac{\qty{5.701e-10}{J}}{\qty{1.602e-19}{J.eV^{-1}}} = \qty{3.558e9}{eV}$.
        Answer: \final{$E = \qty{3.56}{GeV}$} (each beam at least \qty{1.78}{GeV}). The
        electron and positron are far lighter than taus, so this mass comes from the
        kinetic energy the accelerator gave them. Mass--energy is conserved. (The tau was
        found this way at SLAC in 1975.)
  \item Any three: gravity is not part of the Standard Model and the graviton has never
        been found; dark matter and dark energy (about 95\% of the Universe) are not
        explained; we do not know why there is far more matter than antimatter; neutrinos
        have mass, which the first Standard Model did not include; we do not know why there
        are three generations or why particles have the masses they have.
\end{chapteranswers}
